\documentclass[11pt,a4paper]{coverletter}

\senderblock
  {Radheem Bin Razi}
  {Distributed Systems \& Cloud-Native Engineer}
  {Ilmenau, Germany \textbar\ \href{mailto:sheikh.radheem@gmail.com}{sheikh.radheem@gmail.com} \textbar\ \href{https://radheem.github.io/my-cv}{portfolio}}

\begin{document}

%% ===== ENGLISH =====
\recipient{ALTEN Germany}{Hiring Team}{}

\opening{Dear Hiring Team,}

\vspace{8pt}\noindent\textbf{1. Why ALTEN Germany?}\par\vspace{3pt}

Building a simulator that translates raw network traffic into clear, actionable visualizations is the kind of engineering challenge I enjoy most. I am applying for the Python Developer role at ALTEN Germany because your focus on extending communication simulators with OSI layers two through seven aligns directly with my background in protocol-level systems and data-driven tooling.

\vspace{8pt}\noindent\textbf{2. Why me?}\par\vspace{3pt}

My recent work designing a 5G testbed required writing Python xApps that process real-time protocol metrics and route them through a Kafka pipeline. I structured those applications with strict unit tests and version-controlled them in Git, which taught me how to keep simulation code both testable and maintainable. I also built a browser automation workflow that captures and processes external data streams, applying the same disciplined approach to data ingestion and validation. That experience translates directly to extending your existing simulator and adding robust automated testing.

I am comfortable working with networking stacks and have already designed interactive dashboards to surface performance and error patterns for engineering teams. I prefer writing clean, well-documented code that other developers can extend without guessing the original intent. I have also integrated multiple third-party vendors through adapter layers, which reinforced my habit of writing clear API documentation and UML diagrams.

\vspace{8pt}\noindent\textbf{3. Why now?}\par\vspace{3pt}

I am ready to bring this simulation and visualization experience to ALTEN immediately. I am available for full-time work right now and can also accommodate working-student or part-time arrangements if needed. I am currently based in Germany and can relocate to Jena within two to three weeks. My Blue Card status requires no sponsorship overhead, so the employer only needs to issue the contract. I look forward to discussing how I can contribute to your simulator project.

\closing{Sincerely,}{Radheem Bin Razi}

\clearpage
\selectlanguage{ngerman}

%% ===== DEUTSCH =====
\recipient{ALTEN Germany}{Personalabteilung}{}

\opening{Sehr geehrtes Hiring-Team,}

\vspace{8pt}\noindent\textbf{1. Warum ALTEN Germany?}\par\vspace{3pt}

Die Entwicklung eines Simulators, der rohen Netzwerkverkehr in klare, handlungsorientierte Visualisierungen übersetzt, ist genau die Art von ingenieurtechnischer Herausforderung, die ich am meisten schätze. Ich bewerbe mich als Python Developer bei ALTEN Germany, da sich Ihr Fokus auf die Erweiterung von Kommunikationssimulatoren um die OSI-Schichten zwei bis sieben direkt mit meinem Hintergrund in protokollbasierten Systemen und datengetriebener Tooling deckt.

\vspace{8pt}\noindent\textbf{2. Warum ich?}\par\vspace{3pt}

Bei meiner jüngsten Arbeit zur Entwicklung eines 5G testbeds war es erforderlich, Python xApps zu schreiben, die Echtzeit-Protokollmetriken verarbeiten und über eine Kafka pipeline leiten. Ich habe diese Anwendungen mit strengen Unit-Tests strukturiert und in Git versioniert, was mir gezeigt hat, wie Simulationscode sowohl testbar als auch wartbar bleibt. Zudem habe ich einen Browser-Automatisierungsworkflow entwickelt, der externe Datenströme erfasst und verarbeitet, wobei ich denselben disziplinierten Ansatz für Datenerfassung und -validierung angewendet habe. Diese Erfahrung überträgt sich direkt auf die Erweiterung Ihres bestehenden Simulators und die Integration robuster automatisierter Tests.

Ich bin im Umgang mit Netzwerk-Stacks vertraut und habe bereits interaktive Dashboards entwickelt, um Leistungs- und Fehlermuster für Engineering-Teams sichtbar zu machen. Ich bevorzuge es, sauberen, gut dokumentierten Code zu schreiben, den andere Entwickler erweitern können, ohne die ursprüngliche Intention erraten zu müssen. Darüber hinaus habe ich über Adapter-Schichten mehrere Drittanbieter integriert, was meine Gewohnheit gefestigt hat, klare API-Dokumentationen und UML-Diagramme zu erstellen.

\vspace{8pt}\noindent\textbf{3. Warum jetzt?}\par\vspace{3pt}

Ich bin bereit, diese Erfahrung in Simulation und Visualisierung sofort bei ALTEN einzubringen. Ich stehe ab sofort für eine Vollzeitstelle zur Verfügung und kann bei Bedarf auch ein Werkstudenten- oder Teilzeitmodell unterstützen. Ich befinde mich derzeit in Deutschland und kann innerhalb von zwei bis drei Wochen nach Jena umziehen. Mein Blue Card-Status erfordert keine zusätzliche Sponsorship-Abwicklung, sodass der Arbeitgeber lediglich den Arbeitsvertrag ausstellen muss. Ich freue mich darauf, in einem Gespräch zu besprechen, wie ich zu Ihrem Simulator-Projekt beitragen kann.

\closing{Mit freundlichen Grüßen,}{Radheem Bin Razi}

\end{document}
